\subsection{线性方程}

设 E、F 为两个 A-模。形如 \(u(x) = y_{0}\) 的每个方程（其中 \(u: E \to F\) 是一个给定的线性映射， \(y_{0}\) 是F的一个给定元素，而未知量x满足其取值在E中的条件）称为线性方程； \(y_{0}\) 称为方程的右端；如果 \(y_{0} = 0\) ，则方程称为齐次线性方程。

每个元素 \(x_{0} \in E\) ，若它满足 \(u(x_{0}) = y_{0}\) ，则称为线性方程 \(u(x) = y_{0}\) 的解。†

通常粗略地说，如果一个问题的求解等价于确定某个线性方程的解，那么这个问题就是线性的。

给定线性方程 \(u(x) = y_{0}\) ，方程 \(u(x) = 0\) 称为与 \(u(x) = y_{0}\) 相伴的齐次线性方程。

命题14. 若 \(x_{0}\) 是线性方程 \(u(x) = y_{0}\) 的一个解，则该方程的解集等于诸元素 \(x_{0} + z\) 的集合，其中 z 遍历相伴的齐次方程 \(u(x) = 0\) 的解集。

关系 \(u(x) = y_0\) 可以写成 \(u(x) = u(x_0)\) ，它等价于 \(u(x - x_0) = 0\) 。

换句话说，如果方程 \(u(x) = y_{0}\) 至少有一个解 \(x_{0}\) ，其解集就是集合 \(x_{0} + u^{-1}(0)\) ，它是由 u 的核 \(u^{-1}(0)\) 经平移得到的。注意， \(u^{-1}(0)\) 作为一个子模，永远不会是空的，因为它包含 0（称为齐次方程 \(u(x) = 0\) 的零解或平凡解）。

根据命题14，为使线性方程 \(u(x) = y_0\) 恰有一个解，必要且充分的是它至少有一个解，且 \(u^{-1}(0) = \{0\}\) （换句话说，即相伴的齐次方程没有非零解，或者说 \(u\) 是单射）；在这种情况下，对所有 \(y \in F\) ，方程 \(u(x) = y\) 至多有一个解。

命题15. 设 \(u\) 是模 \(\mathbf{E}\) 到模 \(\mathbf{F}\) 的一个线性映射。若方程 \(u(x) = y_0\) 至少有一个解，则 \(y_0\) 正交于 \(^t u\) 的核。

说 \(u(x) = y_{0}\) 有解意味着 \(y_{0} \in u(\mathbf{E})\) ，且该命题由第5号命题8的推论得出。

注意，命题15所给出的、方程 \(u(x) = y_{0}\) 有解的必要判据，当A是域时是充分的（§7，第6号，命题12），但一般不然（习题10）。

注记。(1) 设E是一个A-模， \((\mathbf{F}_{i})_{i\in I}\) 是一族A-模，并且对所有 \(i\in I\) ，设 \(u_{i}:E\to F_{i}\) 是一个线性映射。每个线性方程组

\[
u _ {i} (x) = y _ {i} \quad (i \in I)\tag{29}
\]

其中 \(y_{i} \in \mathbf{F}_{i}\) 是给定的，等价于单个线性方程 \(u(x) = y\) ，其中 \(u\) 是映射 \(x \mapsto (u_{i}(x))\) （从 \(\mathbf{E}\) 到 \(\mathbf{F} = \prod_{i \in I} \mathbf{F}_{i}\) ），而 \(y = (y_{i})\) 。方程组(29)称为齐次的，如果 \(y_{i} = 0\) 对所有 \(i \in I\) 成立。

\begin{enumerate}
\def\labelenumi{(\arabic{enumi})}
\setcounter{enumi}{1}
\tightlist
\item
  假设E有一个基 \((a_{\lambda})_{\lambda \in L}\) ；如果设 \(u(a_{\lambda}) = b_{\lambda}\) 对所有 \(\lambda \in L\) 成立，那么说 \(x = \sum_{\lambda \in L} \xi_{\lambda} a_{\lambda}\) 满足方程 \(u(x) = y_0\) ，就等价于说A中元素的（有限支撑的）族 \((\xi_{\lambda})_{\lambda \in L}\) 满足关系
\end{enumerate}

\[
\sum_ {\lambda \in L} \xi_ {\lambda} b _ {\lambda} = y _ {0}.\tag{30}
\]

反之，寻找A中元素的满足(30)的有限支撑的族 \((\xi_{\lambda})_{\lambda \in L}\) ，等价于求解线性方程 \(u(x) = y_0\) ，其中 \(u\) 是从 E 到 F 的唯一的线性映射，使得 \(u(a_{\lambda}) = b_{\lambda}\) 对所有 \(\lambda \in L\) 成立（§1，第 11 号，命题 17 的推论 3）。

\begin{enumerate}
\def\labelenumi{(\arabic{enumi})}
\setcounter{enumi}{2}
\tightlist
\item
  线性方程 \(u(x) = y_0\) 称为标量的，当 \(F = A_s\) ，从而 \(u\) 是 \(E\) 上的线性形式且 \(y_0\) 是标量时。若 \(E\) 有一组基 \((a_\lambda)_{\lambda \in L}\) ，则由注记（2）可知，这样的方程也可写成
\end{enumerate}

\[
\sum_ {\lambda \in L} \xi_ {\lambda} \alpha_ {\lambda} = y _ {0} \in A\tag{31}
\]

其中标量族 \((\alpha_{\lambda})\) 是任意的，并且约定该族 \((\xi_{\lambda})\) 必须具有有限支撑。一般地，标量线性方程组

\[
\sum_ {\lambda \in L} \xi_ {\lambda} \alpha_ {\lambda i} = \eta_ {i} \quad (i \in I)\tag{32}
\]

其中 \(\alpha_{\lambda i} \in A\) 且 \(\eta_i \in A\) ，其（A中的）解指的是一个有限支撑的族 \((\xi_\lambda)_{\lambda \in L}\) （由 \(A\) 中的元素组成），且满足(32)； \(\alpha_{\lambda i}\) 称为方程组的系数， \(\eta_i\) 称为右端项。这样的方程组的求解等价于方程 \(u(x) = y\) ，其中 \(y = (\eta_i)\) ，而 \(u: A_s^{(L)} \to A_s^I\) 是线性映射

\[
\left(\xi_ {\lambda}\right) \mapsto \left(\sum_ {\lambda \in L} \xi_ {\lambda} \alpha_ {\lambda i}\right).
\]

\begin{enumerate}
\def\labelenumi{(\arabic{enumi})}
\setcounter{enumi}{3}
\tightlist
\item
  线性方程组(32)在需要避免混淆时也称为左标量线性方程组。一个方程组
\end{enumerate}

\[
\sum_ {\lambda \in L} \alpha_ {\lambda i} \xi_ {\lambda} = \eta_ {i} \quad (i \in I)\tag{33}
\]

同样称为右标量线性方程组；把 \(\xi_{\lambda}, \eta_{i}\) 和 \(\alpha_{\lambda i}\) 视为属于A的相反环 A \(^{0}\) ，这样的方程组便可立即转化为方程组(32)。

